% 第 14 章 DRC Waivers
\bichapter{DRC Waivers}{DRC Waivers}
\label{chap:waiver}

本章介绍 IC Validator 工具的 DRC waivers 能力及其用法。

DRC waivers 能力分以下各节说明：
\begin{itemize}
  \item 定义 Waivers
  \item 创建 Waiver 数据库
  \item 使用 Waiver 数据库
\end{itemize}

DRC waivers 允许你对 IC Validator 运行中发现的错误进行分类与标注，并将错误导出到 waiver 数据库，供后续 IC Validator 运行使用。
DRC Waivers 流程分为三个独立步骤：
\begin{itemize}
  \item 定义哪些错误可被 waive，并通过 VUE 或 \cmd{icv\_pydb} 命令行实用程序对错误数据库加注
  \item 在磁盘上创建持久化 waiver 数据库
  \item 在后续 IC Validator 运行中应用 waivers
\end{itemize}

以下各节分别说明这些步骤。

\figplaceholder{Figure 96: General Waiver Flow}{一般 Waiver 流程}{fig:waiver-flow}

% ============================================================
\section{定义 Waivers}
\subsection*{Defining Waivers}

可使用 VUE 或 ASCII 错误报告文件 \cmd{cell.LAYOUT\_ERRORS} 审查写入错误数据库（PYDB）的设计错误。
在许多设计中，存在可被 waive 的 DRC 违例，即这些错误无需在版图中修复。
这可以是必须在 tapeout 前修复的临时 waiver，也可以是晶圆厂批准的永久 waiver。
可为 waivers 添加用户注释或预定义的 waiver 属性，以便为审查 waivers 的其他工程师提供更多上下文。

IC Validator SQLite 错误数据库允许多用户并发访问；在审查错误并应用 waivers 时，访问该数据库的其他用户可以看到你已应用的更新。

IC Validator 支持若干 waiver 类别：
\begin{itemize}
  \item 持久化 waiver 类别：Waive、Watch、Ignore
  \item 非持久化 waiver 类别：Error、Fixed
\end{itemize}

持久化 waiver 类别可导出到 waiver 数据库。见``创建 Waiver 数据库''。
通过多种持久化 waiver 类别，可将已 waive 的错误划分到不同类别，使设计过程中的 waiver 管理更简便。
这些 waiver 类别同等对待。例如，可在整个设计过程直至 tapeout 均拟 waive 该错误时选择 Waive，并选择 Watch 以标示可在下一两次设计迭代中修复的错误。

非持久化 waiver 类别向调试同一错误数据库的其他用户提供信息。
例如，将错误定义为 Error 并附用户注释 ``this error needs to be fixed in conjunction with rule ABCD''，则在同一错误数据库中审查该错误的任何人都能看到此指导。

\subsection{使用 VUE 定义 Waivers}
\subsubsection*{Using VUE to Define Waivers}

定义 waivers 最常见的方式是通过 VUE DRC 浏览器。有多种方式 waive 错误。

要 waive 某违例（规则）中的全部错误形状，在 Violation Browser 中选择该违例并右键单击以查看选项菜单。选择适当的 waiver 类别，如图~\ref{fig:vue-viol-browser} 所示。

\figplaceholder{Figure 97: Violation/Cell in Violation Browser}{Violation Browser 中的违例/单元}{fig:vue-viol-browser}

本例选择了 Waive 类别。选择后出现窗口，允许提供可选用户注释并定义 waiver 上下文。

\figplaceholder{Figure 98: Optional User Comment}{可选用户注释}{fig:vue-user-comment}

Context 选项决定这些 waivers 在未来 IC Validator 运行中何时适用。设为 ALL 时，若与错误匹配则始终应用该 waiver。设为顶层单元时（本例顶层单元名为 ``top''），仅当当前顶层 block 再次作为顶层单元运行时才应用该 waiver。

定义 waivers 后，错误浏览器中会新增一列以显示含已 waive 错误的违例。

\figplaceholder{Figure 99: Error Browser Column}{错误浏览器列}{fig:vue-err-browser-col}

每个错误现已用 waiver 信息加注。

\figplaceholder{Figure 100: Waiver Information}{Waiver 信息}{fig:vue-waiver-info}

状态现为 Waive。From 列显示 ``vue''，表示当前 waiver 由用户应用，并非 IC Validator 引用 waiver 数据库时找到的。

Error List 可用于仅 waive 某些错误，如图~\ref{fig:vue-error-list} 所示。

\figplaceholder{Figure 101: Error List}{错误列表}{fig:vue-error-list}

可通过按住 Shift-click 与 Ctrl-click 选择 R.M4.W 违例中的个别错误，然后右键单击查看 waive 选项菜单。另有选项 Waive In Area，允许 waive 特定区域内的全部错误，而无需从 Error List 中逐个选择。

错误被 waive 后，可在错误列表中过滤以隐藏特定 waiver 类别。在 DRC/Extraction VUE 选项中，使用 Enable/Disable filter functions 部分选择要在 Error List 中显示的 waiver 类别。

\figplaceholder{Figure 102: Filter Settings}{过滤设置}{fig:vue-filter}

本例中，分类为 Waive 的错误从 Error List 中过滤掉。显示所有其他 waiver 类别。

图~\ref{fig:vue-updated-list} 展示更新后的 Error List：

\figplaceholder{Figure 103: Updated Error List}{更新后的错误列表}{fig:vue-updated-list}

\subsection{使用命令行实用程序定义 Waivers}
\subsubsection*{Defining Waivers Using the Command-Line Utilities}

可使用 \cmd{icv\_pydb} 命令行实用程序定义 waivers。更多信息见第~\ref{chap:errmgmt}~章``错误管理实用程序''。

\subsection{自定义 Waiver 属性}
\subsubsection*{Custom Waiver Properties}

除用户注释外，还可为给定规则与单元定义额外的 waiver 属性。这些属性有助于添加信息，并提供跨设计或工艺节点管理 waivers 的额外方式。
例如，可为 waivers 赋予定义特定跟踪号及其当前管理审批状态的 waiver 属性。若 waiver 未获批准，可使用 \cmd{pydb\_export} 与 \cmd{-select\_waiver\_properties} 命令行选项，从 waiver 数据库中移除具有该跟踪号的任何 waiver。

允许的值定义在由 VUE 读取的 JSON 格式文件中。

\begin{lstlisting}[style=shell,caption={JSON 属性文件（waiver\_properties.json）},label={ex:waiver-props-json}]
{
    "waiver_property_definitions": [
      {
         "property_name": "WaiverTracking#",
         "property_values": []
      },
      {
         "property_name": "Valid_through",
         "property_values": ["milestone1", "milestone2", "milestone3"]
      },
      {
         "property_name": "Status",
         "property_values": ["pending", "approved", "not approved"],
         "default_value": "pending"
      }
    ]
}
\end{lstlisting}

每个属性有 \cmd{property\_name}、\cmd{property\_values}，以及可选的 \cmd{default\_value}。
\cmd{property\_name} 遵循 OASIS 属性标准，即使用可打印 ASCII 字符，排除空格、制表符、换行符与控制字符。属性值不受限制。

\begin{noteBox}
\cmd{WAIVER\_REPORT} 文件使用 \texttt{|} 作为多个属性值之间的分隔符。为避免混淆，请勿使用该字符。
\end{noteBox}

设置下列环境变量后，VUE 会读取该 JSON 文件：
\begin{lstlisting}[style=shell]
setenv ICV_WAIVER_PROPERTY_DEFINITION_FILE waiver_properties.json
\end{lstlisting}

然后创建含允许属性的自定义窗口。
在 Violation Browser 中，右键单击选择规则与单元。此时可用 Add Waiver Properties 选项。

\figplaceholder{Figure 104: Add Waiver Properties}{添加 Waiver 属性}{fig:vue-add-waiver-props}

选择该选项以定义用户属性。

该规则/单元的所有已 waive 错误继承这些属性。可在 Violation Detail 窗格中查看属性，如下所示。

% ============================================================
\section{创建 Waiver 数据库}
\subsection*{Creating a Waiver Database}

为当前设计定义全部 waivers 后，将信息导出到独立的 waiver 数据库。Waiver 数据库可以是 SQLite 格式（CPYDB）或 OASIS 格式。Waiver 数据库是自包含的，即 waiver 错误所需的全部信息（用户注释、用户 ID、日期、waiver 属性等）均包含在数据库内。

可使用 VUE 或 \cmd{pydb\_export} 命令行实用程序导出 waiver 数据库。

要用 VUE 导出 SQLite waiver 数据库，选择 Classification > Export Classification Data (CPYDB)。

在此可定义数据库的详细信息。

输出数据库名写入启动 VUE 时的工作目录。要写到其他路径，在输出文件名中包含该路径。默认情况下，IC Validator 创建新的 waiver 数据库。可通过设置 Append Mode 将此信息追加到现有 CPYDB 数据库。导出期间，VUE 导出错误数据库中定义的任何新 waivers。要包含全部 waivers（包括先前从另一 waiver 数据库匹配到的），勾选 Include Matched Errors 框。

SQLite 格式允许为 waiver 数据库设置口令。除非知道口令，否则无法修改该 waiver 数据库。

要从 VUE 导出 OASIS 格式 waiver 数据库，选择 Classification > Export Classification Data (OASIS)。

此处窗口提供 waiver 数据库的设置。

输出数据库名写入启动 VUE 时的工作目录。要写到其他路径，在输出文件名中包含该路径。输出的 OASIS 格式将 waiver 形状放在特殊命名的单元中，该单元为每组 waivers 提供规则与单元名（\cmd{Rule\_Name@cellname}）。Rule Name Delimiter 允许你定义工艺节点规则名的分隔符。在下列示例中，规则名使用 ``:'' 字符作为规则名与规则描述之间的分隔符。设置后，分隔符左侧的全部字符定义 waiver 单元名。若未设置，或规则名很长，则 Max Rule Name Length 将该 waiver 单元名截断到该长度。默认值为 \cmd{"[[:space:]]+:|[[:space:]]*:[[:space:]]"}，可匹配大多数晶圆厂 runset 命名约定。

现可在 IC Validator WorkBench 中打开 OASIS 格式 waiver 数据库。在 waiver 数据库的树视图中（如下图所示），可以看到数据库有一个以 \cmd{@} 为前缀的单元。这是 holding cell，包含为该单元定义的全部 waivers。该单元本身不含任何 waiver 形状；形状包含在放置于 holding cell 内的规则专用单元中。

每个 waiver 单元（\cmd{rule@cell}）在单元上附着属性。

这些属性在导出时设置。导出 waiver 数据库的用户会将其用户 ID 列在 \cmd{createID} 属性中，以及导出时间戳。\cmd{waiverChecksum} 属性在导出时对该单元计算。IC Validator 读取 waiver 单元时，会计算校验和并与存储的属性比较。若属性不同，表示数据库自导出后已被修改。\cmd{WAIVER\_REPORT} 文件末尾包含校验和是否匹配的信息。

给定 Violation/Cell 的全部 waiver 形状写入任意层号。IC Validator 不使用层号，而是使用 \cmd{violationComment} 属性来了解 waiver 形状适用于哪条规则。因此不需要映射文件。

% ============================================================
\section{使用 Waiver 数据库}
\subsection*{Using a Waiver Database}

创建 waiver 数据库后，可将其用作其他 IC Validator 运行的输入。新运行不必是同一顶层单元。IC Validator 导入当前运行中存在的 Cell / Violation 的全部 waiver 形状。IC Validator waivers 是完全层次化的。

\cmd{waiver\_options()} 函数提供对 IC Validator 运行期间如何应用与报告 waivers 的控制。该命令分为两个主要部分。第一部分是数据库管理：提供 waiver 数据库列表并定制每个数据库的处理方式。第二部分提供 IC Validator 对全部 waivers 与报告使用的全局 waiver 控制。

数据库管理由 \cmd{databases} 参数控制。IC Validator 自动检测数据库格式，同一 IC Validator 运行中可同时提供 CPYDB 与 OASIS 格式。使用 OASIS 格式 waiver 数据库时不需要层映射文件。

对每个 waiver 数据库，可定制控制以提供跨设计的灵活性。包括：
\begin{itemize}
  \item 控制在找不到 waiver 数据库时如何处理。
  \item 用此数据库 waive 的错误是抑制还是分类。抑制意味着已 waive 的错误不写入 PYDB 或 \cmd{LAYOUT\_ERRORS} 文件。
  \item 也可为 waiver 数据库给出单元前缀或单元后缀。若流程在将单元放入更大宏块时使单元唯一化，但基单元的 waivers 仍可应用于重命名后的单元，则此功能很有用。
  \item 生成特殊报告，列出 IC Validator 运行期间未匹配的数据库中全部 waiver 形状。
\end{itemize}

全局 waiver 控制决定如何将 waiver 形状与运行期间发现的错误形状比较，包括任何规则专用覆盖与报告生成。

\subsection{匹配模式}
\subsubsection*{Matching Modes}

可控制 IC Validator 如何将 waiver 形状匹配到给定错误形状。所有匹配均在所有相交的 waiver 形状已合并、所有相交的错误形状已合并之后进行。在某些情况下，为允许合并发生，形状可能需要提升到共同父单元。匹配是层次化的，即层次合并后的 waiver 形状不必与层次合并后的错误形状在同一单元中即可找到匹配。

有三种匹配模式可用：\cmd{EXACT}、\cmd{ENCLOSE}（默认）与 \cmd{INTERACT}。
\begin{itemize}
  \item \cmd{EXACT} 模式：层次合并后的 waiver 形状必须与层次合并后的错误形状完全匹配，才会应用 waiver。
  \item \cmd{ENCLOSE} 模式是默认匹配模式。Waiver 形状必须包围错误形状。
  \item \cmd{INTERACT} 模式：任何与 waiver 形状层次相交的错误均被 waive。使用 \cmd{include\_touch} 参数可将外侧相触（outside-touch）交互视为 waive 的相交。
\end{itemize}

\begin{noteBox}
使用 \cmd{tolerance} 参数可改善全角度（all-angle）形状的匹配。
\end{noteBox}

在 \cmd{waiver\_options()} 中设置的匹配模式适用于全部 \cmd{waiver\_databases}。可用 \cmd{rule\_overrides} 参数覆盖特定规则的匹配模式。例如，除 \cmd{R.MY.STRICT\_RULE} 外的全部规则允许 \cmd{ENCLOSE} 模式；对该规则，仅当 waiver 与错误形状完全匹配时才应用 waivers。例如：
\begin{lstlisting}[style=pxl]
waiver_options(
  databases = { {"/path/to/waiver.oas"}},
  matching_mode = ENCLOSE,
  rule_overrides = {
    violation_comments = {"R.MY.STRICT_RULE"},
    matching_mode = EXACT,
    tolerance     = 0.005
  }
 );
\end{lstlisting}

\subsection{Waiver 报告文件}
\subsubsection*{Waiver Report File}

IC Validator 将全部已分类（已 waive）的错误写入 IC Validator 工作目录中的 waiver 报告文件（\cmd{cell.WAIVER\_REPORT}）。

\cmd{LAYOUT\_ERRORS} 文件与 VUE run\_summary 选项卡显示：
\begin{lstlisting}[style=shell]
ERRORS (WAIVERS APPLIED)
\end{lstlisting}

默认生成 waiver 报告文件，含规则汇总（包括已应用的 waivers）以及输入 waiver 数据库列表。
\begin{lstlisting}[style=shell]
-------------------------------------------------------------------------
----
WAIVER SUMMARY
-------------------------------------------------------------------------
----

R.M3.L: M3 can not have 3 or mor... | Waive : 5 | Ignore : 0 | Watch : 0
 | WaiverConflict : 0 | Unmatched : 0
R.M4.W: M4 Minimum width >= 0.24    | Waive : 4 | Ignore : 0 | Watch : 0
 | WaiverConflict : 0 | Unmatched : 0
\end{lstlisting}

表列为：规则名、Waive、Ignore、Watch 类别计数、WaiverConflict 计数，以及 Unmatched waiver 计数。
\begin{itemize}
  \item 当具有不同 waiver 类别的相交 waiver 形状被层次合并在一起时，发生 waiver 冲突。
  \item Unmatched 是该规则在所有输入 waiver 数据库中未在当前设计中匹配的 waiver 形状计数。
\end{itemize}

也可通过设置下列内容生成详细 waiver 报告：
\begin{lstlisting}[style=pxl]
waiver_options(report_waiver_details = true)
\end{lstlisting}

这会为每个 waiver 添加表格，含单元、waiver 类别、匹配模式、上下文、匹配到的 waiver 形状所在的 waiver 数据库、任何 waiver 属性，以及错误详情（距离、包围盒等）。

\begin{noteBox}
可用 \cmd{-runset\_config} 命令行选项修改 \cmd{waiver\_options()} 函数。
\end{noteBox}

可在 IC Validator 运行完成后生成 waiver 报告文件。若不想在默认 IC Validator 运行中查看详情，但希望运行后再查看，此功能很有用。
\begin{lstlisting}[style=shell]
%icv_pydb -i ./run_details/pydb/PYDB_top \
     -create_waiver_report top.DETAILED_WAIVERS \
     -include_waiver_details
\end{lstlisting}

也可用 \cmd{icv\_pydb} 实用程序生成 JSON 格式的 waiver 报告文件。若需编写脚本处理给定设计上已应用的 waivers，这种结构化格式更易于解析。
\begin{lstlisting}[style=shell]
% icv_pydb -i ./run_details/pydb/PYDB_top \
      -create_json_waiver_file top.waiver_file \
      -include_waiver_details
\end{lstlisting}

\begin{lstlisting}[style=shell,caption={JSON 格式 Waiver 报告文件},label={ex:json-waiver-report}]
{
     "icv_run_info": {
     "library_name": "../../../lab_data/INLIB_TRAINING.GDS",
     "structure_name": "top",
     "generated_by": "IC Validator RHEL64
  S-2021.06-SP3-BETA-20211123.7086978 2021/11/23",
     "runset_name": "../modular.rs",
      "user_name": "johndoe",
      "time_started": "2021/11/24 12:00:58PM",
      "time_ended": "2021/11/24 12:01:13PM"
   },
   "rules": [
      {"rule": "R.M3.L", "violation_comments": ["R.M3.L: M3 can not have 3
 or more consecutive line segments be < 0.07000000000000001"]},
      {"rule": "R.M4.W", "violation_comments": ["R.M4.W: M4 Minimum width
 >= 0.24"]}
   ],
   "waiver_databases": [
      {
        "waiver_database_id": "WDB1",
        "file":
 "/path/to/johndoe/TRAINING/drc/lab_answers/lab5_modularize/mywaivers_db.oas",
        "checksum": "PASS"
      }
   ],
   "waiver_details": [
      {
        "rule": "R.M3.L",
        "command": ".././include/metal_rules.rs:78:vertices_edge",
        "cell": "top",
        "class": "Waive",
        "mode": "Enclose",
        "context": "ALL",
        "waiver_sources": ["WDB1:top"],
        "unique_id": "E.12.100.1",
        "bbox": [30.2400, 15.3700, 30.2400, 15.4000],
        "user_comments": ["waiving M3.L rule."]
      }
],
   "unmatched_waiver_details": [ ],
   "rule_overrides": [ ]
}
\end{lstlisting}
